\chapter*{Abstract}
% \phantomsection
% \addcontentsline{toc}{chapter}{Abstract}

Two-dimensional light-element xenes provide a compact platform for exploring quantum properties generated by reduced dimensionality, strong covalent bonding, and low atomic mass.
Graphene established the basic route toward atomically thin quantum materials, but its zero band gap limits several electronic and optoelectronic applications.
This motivates the search for alternative elemental and compound xenes in which bonding topology, stacking, and surface functionalization can tune the electronic structure and optical response.
In this thesis, we investigate boron- and beryllium-based two-dimensional systems using first-principles calculations based on density functional theory.
The calculations determine atomic structure, electronic band structure, density of states, dynamical stability, dielectric response, and derived optical properties.
For selected systems, the analysis also characterizes magnetism, superconductivity, and non-trivial topology.

The graphene-based heterostructure study examines van der Waals bilayers formed from graphene, borophene, and two-dimensional boron carbides, namely \ce{BC3} and \ce{B4C3}.
The isolated monolayers provide contrasting electronic characters.
Graphene remains semimetallic, borophene is metallic, whereas monolayer \ce{BC3} and monolayer \ce{B4C3} are indirect-gap semiconductors with HSE06 band gaps of \qty{1.822}{\electronvolt} and \qty{2.381}{\electronvolt}, respectively.
When these materials form bilayer heterostructures, Graphene-Borophene remains metallic, whereas Graphene-\ce{BC3} and Graphene-\ce{B4C3} are semimetallic.
The graphene-derived linear bands remain near the $\mathrm{K}$ point, with small band-gap openings of \qty{20}{\milli\electronvolt} to \qty{50}{\milli\electronvolt}.
The calculated dielectric response and derived optical properties show visible-region absorption and plasmonic excitations in all three heterostructures.
Graphene-Borophene gives the strongest absorption and supports low-energy plasmonic features.
Graphene-\ce{B4C3} also displays non-negligible off-diagonal components in the dielectric tensor $\varepsilon_{\alpha\beta}(\omega)$, indicating anisotropic and possibly chiral optical behaviour.

The penta-bipyramid boron study examines boron phases derived from bulk orthorhombic \ce{o-B14}.
The edge-sharing pentagonal bipyramids create an unusual bonding network and give boron access to several quantum phases within one structural family.
The calculations investigate bulk \ce{o-B14}, monolayer \ce{o-B14}, bilayer \ce{o-B14}, and the hydrogen-terminated bilayer.
The two-dimensional forms display diverse electronic behaviour.
The bilayer is semiconducting with an HSE06 band gap of \qty{0.67}{\electronvolt}.
The monolayer shows magnetism, whereas the hydrogen-terminated bilayer becomes superconducting with an anisotropic Migdal--Eliashberg critical temperature $T_\mathrm{c}$ of \qty{24.3}{\kelvin}.
The optical response is strongly anisotropic.
Several systems absorb in the visible region, and the metallic phases support directional plasmonic modes.
These results demonstrate that the electron-deficient bonding of boron can support semiconductivity, magnetism, superconductivity, and anisotropic optical response in closely related structures.

The beryllene study investigates pristine and hydrogen-functionalized two-dimensional beryllene phases.
Beryllium represents an extreme light-element limit for elemental two-dimensional materials.
This thesis considers pristine $\alpha$-beryllene, $\beta$-beryllene, a cubic trilayer beryllene phase, and their hydrogen-functionalized counterparts.
The stability analysis confirms the stability of the previously proposed $\alpha$- and $\beta$-beryllene phases and identifies a dynamically stable cubic trilayer beryllene phase.
Hydrogen functionalization strongly modifies the electronic structure.
It drives $\alpha$-beryllene into a wide-gap semiconducting state, whereas the hydrogen-functionalized $\beta$-based structures retain metallic behaviour.
The dielectric response reveals strong anisotropy, visible-to-ultraviolet absorption, modified plasmonic behaviour, and direction-dependent screening.
Using parity analysis of spin-orbit-coupled electronic states, the calculations further identify non-trivial $\mathbb{Z}_2$ topology in pristine $\alpha$-beryllene and cubic trilayer beryllene, while pristine $\beta$-beryllene remains topologically trivial.

These studies demonstrate that light-element xenes can host complex quantum and optical behaviour when their atomic structures are organized in low-dimensional forms.
Heterostructure formation, bonding topology, reduced dimensionality, and hydrogen functionalization provide practical routes for tuning electronic structure, dielectric response, plasmonic modes, superconductivity, magnetism, and topology.
The results of this thesis therefore provide specific predictions for boron and beryllium allotropes, and outline a broader strategy for designing ultrathin quantum materials from light elements.
